\par\addvspace{12pt}\Needspace{9\baselineskip}

\section{独立性}

% ITEM: probability-example-013
\solutionitem{probability-example-013}{例6.13}
（据2022年全国乙卷理科改编）
某棋手与甲、乙、丙各比赛一盘，各盘结果相互独立，获胜概率分别为\(0<p_1<p_2<p_3<1\)。记三盘中至少出现一段连续两盘获胜的概率为\(q\)。比较三种不同的第二盘对手，确定怎样安排使\(q\)最大。
\solutionanswer
记按时间排列的获胜概率为\(a,b,c\)。连续两胜即前两盘皆胜或后两盘皆胜，由容斥得
\[
 q=ab+bc-abc=(ab+bc+ca-abc)-ac.
\]
调换次序，括号中的量不变，只需比较首末两盘的概率之积。因\(p_1p_2<p_1p_3<p_2p_3\)，第二盘对手分别为丙、乙、甲时，有\(q_{\text{丙}}>q_{\text{乙}}>q_{\text{甲}}\)。故丙排第二最优，甲、乙在首末的次序不影响结果。


% ITEM: probability-exercise-06-01
\solutionitem{probability-exercise-06-01}{习题6.6（1）}
独立重复试验5次，每次成功概率为\(0.4\)。求全部成功、至少成功一次及恰成功两次的概率。
\solutionanswer
全部成功\(0.4^5=0.01024\)；至少一次\(1-0.6^5=0.92224\)；恰两次\(\binom520.4^2\,0.6^3=0.3456\)。最后一式先指定成功位置再相加。


% ITEM: probability-exercise-06-02
\solutionitem{probability-exercise-06-02}{习题6.6（2）}
两个独立元件正常工作的概率分别为\(0.9,0.8\)。串联系统要求二者都正常，并联系统要求至少一个正常。求两种系统正常的概率。
\solutionanswer
串联要求同时正常，概率\(0.9\times0.8=0.72\)；并联失败当且仅当二者都失败，故正常概率\(1-0.1\times0.2=0.98\)。


% ITEM: probability-exercise-06-03
\solutionitem{probability-exercise-06-03}{习题6.6（3）}
已知\(P(A)=P(B)=1/2\)，\(A,B\)独立。求\(P(A\cup B)\)及\(P(A\mid A\cup B)\)。观察获得新信息后，所用模型发生了什么变化。
\solutionanswer
交概率为\(1/4\)，故并概率\(3/4\)。由于\(A\subset A\cup B\)，\(P(A\mid A\cup B)=P(A)/P(A\cup B)=2/3\)。新条件删去“两者均不发生”的结果并重新归一化；此时两事件的条件概率均为\(2/3\)，同时发生为\(1/3\)，不再独立。


% ITEM: probability-exercise-06-04
\solutionitem{probability-exercise-06-04}{习题6.6（4）}
独立公平抛三次硬币，令\(A\)表示前两次相同，\(B\)表示后两次相同，\(C\)表示第一、三次相同。判断三事件是否两两独立、相互独立。
\solutionanswer
每个相等事件概率\(1/2\)。任意两个同时发生均等价于三次全同，概率\(2/8=1/4\)，故两两独立。但三者交概率仍为\(1/4\ne(1/2)^3\)，不是相互独立。


% ITEM: probability-exercise-06-05
\solutionitem{probability-exercise-06-05}{习题6.6（5）}
（据2021年天津卷改编）甲、乙每次各猜一个谜语，猜对概率分别为\(5/6,3/5\)。每次两人的结果互不影响，各次活动也互不影响。一人猜对、另一人猜错时，猜对者获胜，否则平局。求甲一次获胜及三次活动中至少获胜两次的概率。
\solutionanswer
甲一次获胜概率\((5/6)(1-3/5)=1/3\)。三次活动独立，至少获胜两次的概率\(\binom32(1/3)^2(2/3)+(1/3)^3=7/27\)。


% ITEM: probability-exercise-07-01
\solutionitem{probability-exercise-07-01}{习题6.7 *}
本题先等可能选机器甲或乙，再用所选机器独立制造两件产品；甲、乙不合格率分别为\(0.1,0.9\)。

在上述两台机器的模型中，已知第一件不合格，求第二件也不合格的概率；再求已知第一件合格时第二件不合格的概率。比较两者与事先的\(0.5\)。
\solutionanswer
两件均不合格的概率\(0.41\)，第一件不合格概率\(0.5\)，故\(P(D_2\mid D_1)=0.82\)。第一合格、第二不合格的联合概率为\(\frac12(0.9\times0.1)+\frac12(0.1\times0.9)=0.09\)，故\(P(D_2\mid D_1^c)=0.18\)。它们分别高于和低于原概率\(0.5\)：第一件结果改变了对机器来源的判断。
\par
